\section{Відкладені проєктні рішення та виключення версії 1}\label{sec:vidkladeni}

У цьому розділі зафіксовано свідомі обмеження поточної специфікації: інженерні рішення, вибір яких відкладено на етап проєктування реалізації, а також функції, які явно виключено з першої версії системи.

\subsection{Відкладені рішення щодо технологічного стеку та фреймворків}\label{sec:vidkladeni-stek}
Конкретна мова програмування, вебфреймворк для побудови API, система керування базами даних (СКБД) та бібліотеки для автоматизованого тестування у цій специфікації навмисно не фіксуються~\sked{K42,Q5}. Їх вибір здійснюється розробником на наступному етапі проєктування реалізації та погоджується перед початком написання коду~\sked{K15,K42,Q5}.

При цьому до обраного технологічного стеку встановлюються чіткі обов'язкові критерії (\tabref{tab:tech-criteria})~\sked{K43,Q5}.

\begin{table}[htbp]
\caption{Обов'язкові критерії до вибору технологій реалізації}\label{tab:tech-criteria}
\centering
\begin{tabularx}{\textwidth}{@{}l L@{}}
\toprule
\textbf{Вимога / Критерій} & \textbf{Технічне обґрунтування та значення для системи} \\
\midrule
Підтримка TDD & Наявність зрілих інструментів для юніт- та інтеграційного тестування з високою швидкістю виконання тест-сьютів. \\
Рольове розмежування & Ефективні вбудовані засоби автентифікації та перевірки прав доступу на рівні бекенду (RBAC). \\
Транзакційна атомарність & Підтримка транзакційного механізму в базі даних (ACID / атомарні операції) для гарантії захисту від подвійних бронювань. \\
Персистентність даних & Надійне збереження даних на диску з можливістю коректного відновлення після рестарту. \\
\bottomrule
\end{tabularx}
\end{table}

\subsection{Погоджені виключення та обмеження першої версії}\label{sec:vidkladeni-vykliuchennia}
Відповідно до рішень автора проєкту, у версії 1 не реалізуються наступні функціональні можливості (\tabref{tab:v1-exclusions})~\sked{K25,K28,K32,K33,K34,K49}.

\begin{table}[htbp]
\caption{Погоджені функціональні виключення першої версії}\label{tab:v1-exclusions}
\centering
\begin{tabularx}{\textwidth}{@{}l L L@{}}
\toprule
\textbf{Функціональна область} & \textbf{Виключення у версії 1} & \textbf{Погоджений статус v1} \\
\midrule
Автентифікація & Гостьове бронювання та сторонній вхід (Google, OAuth, SSO). & Лише реєстрація через email/пароль. \\
Оплата квитків & Онлайн-оплата, платіжні шлюзи, статуси оплат, повернення коштів. & Безоплатне резервування (оплата в касі). \\
Ціноутворення & Збереження, розрахунок та відображення цін квитків. & Ціни не відображаються у системі. \\
Бронювання & Тимчасове утримання місць (hold) із таймаутом на час вибору. & Атомарна фіксація в момент кліку. \\
Адміністрування & Керування кількома адміністраторами та правами адмінів. & Єдиний системний адміністратор. \\
Зміна сеансів & Автоматичне перенесення бронювань та email/SMS розсилки. & Ручне перебронювання користувачами. \\
\bottomrule
\end{tabularx}
\end{table}
